\documentclass[11pt]{article}
\usepackage[margin=25mm]{geometry}
\usepackage{fontspec}
\setmainfont{Latin Modern Roman}
\usepackage{amsmath,amssymb,mathtools}
\usepackage{unicode-math}
\setmathfont{Latin Modern Math}
\usepackage{xurl,hyperref}
\hypersetup{colorlinks=true,urlcolor=blue}
\setcounter{secnumdepth}{0}
\setlength{\parindent}{0pt}
\setlength{\parskip}{5pt}
\setlength{\emergencystretch}{3em}
\newcommand{\sref}[2]{\hyperlink{source:#1:#2}{[S#2]}}
\newcommand{\cataloguescope}[1]{\paragraph{Recorded target.}#1}
\newcommand{\R}{\mathbb R}
\newcommand{\C}{\mathbb C}
\newcommand{\N}{\mathbb N}
\newcommand{\Z}{\mathbb Z}
\newcommand{\Q}{\mathbb Q}
\newcommand{\F}{\mathbb F}
\newcommand{\eps}{\varepsilon}
\begin{document}
% UnsolvedMath ID 30000179; source catalogue code OWR-785-005
\setcounter{section}{30000178}
\section{Connections Between CGLMP Inequalities and Quantum Field Theory}\label{30000179}
{\small UnsolvedMath ID 30000179 · Mathematical Physics}
\subsection{Definitions and mathematical statement}

\paragraph{Shared notation.}$\mathbb N=\{1,2,\ldots\}$; the standard number fields have their usual meanings. All additional conventions are stated locally.


\paragraph{Definitions and conventions.} For two parties with settings $1,2$ and outcomes in $\{0,\ldots,d-1\}$, $d\ge2$, the Collins--Gisin--Linden--Massar--Popescu (CGLMP) inequality is a Bell inequality constraining local hidden-variable correlations. In the report it is written using expectations of outcome differences reduced modulo $d$. The quantum Fourier transform on $\mathbb C^d$ is the unitary $F_d|j\rangle=d^{-1/2}\sum_{k=0}^{d-1}e^{2\pi ijk/d}|k\rangle$. Fourier-transform measurement bases enter the report's numerically optimized high-dimensional Bell experiments.

\paragraph{Source scope and qualifications.} The catalogue title and clean statement expand QFT as quantum field theory. The primary report explicitly discusses the Quantum Fourier Transform immediately before asking for the CGLMP--QFT link. Thus the source-context question concerns the Fourier transform; no quantum-field-theory connection is asserted here. The catalogue wording is preserved below to make this discrepancy visible.

\paragraph{Primary-source formulation.} Explain the mathematical connection between the CGLMP inequality and the quantum Fourier transform appearing in the report's optimal measurement constructions.

\paragraph{Catalogue mathematical target.}
\paragraph{Catalogue wording retained for traceability.}
What is the mathematical link between the CGLMP inequality and quantum field theory?

\paragraph{Saved literature context.} The supplied catalogue's assessment dated 2026-08-21 records: The CGLMP inequality has extensive subsequent study, including its optimal violations and information-theoretic interpretations, but no generally accepted mathematical link to quantum field theory resolving the stated problem was found. The assessment is retained as dated provenance; it is not a new comprehensive literature or open-status certification.
\subsection{Short English statement}
Explain the mathematical link between CGLMP Bell inequalities and the quantum Fourier transform used in the optimal measurements.
\subsection{Research attempt}
\paragraph{Model, date and investigated scope.}
OpenAI GPT-6 Astra (configured model \texttt{gpt-6-astra}), ultra reasoning; generated 11 October 2026. Derives modular-difference probabilities for Fourier-phase measurements, honoring the source correction from field theory to Fourier transform.

\paragraph{Sources and literature check.}
The full supplied definition and its explicit QFT correction were reviewed. The Collins et al. arXiv abstract [S4] was opened; no theorem of global optimality was taken from it. The calculation is an elementary reconstruction, with novelty and current status unchecked.

\paragraph{Context and investigated scope.} The full supplied definition and its explicit QFT correction were reviewed. The Collins et al. arXiv abstract [S4] was opened; no theorem of global optimality was taken from it. The calculation is an elementary reconstruction, with novelty and current status unchecked.

\paragraph{The source-corrected Fourier question.} Work with the primary-source meaning of QFT, the quantum Fourier transform; the retained catalogue phrase ``quantum field theory'' is not the target of this calculation. Let $d\geq2$, $\sum_{j=0}^{d-1}|c_j|^2=1$, and $|\psi\rangle=\sum_jc_j|j,j\rangle$. For real phases $\alpha_x,\beta_y$, define orthonormal measurement vectors
\[
 |a_x\rangle=d^{-1/2}\sum_j e^{2\pi i j(a+\alpha_x)/d}|j\rangle,
 \quad |b_y\rangle=d^{-1/2}\sum_j e^{2\pi i j(-b+\beta_y)/d}|j\rangle.
\]
Orthogonality follows from the finite geometric sum for distinct integer outcomes. Their probability is
\[
 p(a,b|x,y)=d^{-2}\left|\sum_jc_j e^{-2\pi i j(a-b+\alpha_x+\beta_y)/d}\right|^2.
\]
There are exactly $d$ pairs with a fixed residue $a-b=k$, so the difference probability is $d^{-1}$ times the squared modulus in this formula with $a-b=k$.

\paragraph{Concrete connection and its limit.} Every CGLMP functional written as a weighted sum of modular-difference probabilities becomes a quadratic form $c^*Hc$ after these substitutions. Explicitly a term with weight $w$ and phase offset $t$ contributes $H_{j\ell}=(w/d)e^{2\pi i(j-\ell)t/d}$. Thus Fourier phases convert cyclic outcome differences into interference between Schmidt coefficients. This is a direct explanation of why the transform is compatible with the inequality's cyclic structure.

The calculation holds also for nonuniform $c_j$ and arbitrary real offsets; at integer offsets and uniform $c_j$, the geometric sum enforces a sharp residue, a useful normalization check. It does not show that these bases maximize the Bell functional over all measurements and states. Proving that reduction, or identifying another precise interpretation of the source's requested ``link'', is the first gap. No quantum-field-theory connection is inferred.

\paragraph{Verification and adversarial audit.}
Checked basis orthogonality, factor d in residue probabilities, quadratic-form signs, and the catalogue/source QFT discrepancy. Model self-audit only; no independent human proof checking. Prompts A and B were applied to this canonical record and its qualified source scope. Checked basis orthogonality, factor d in residue probabilities, quadratic-form signs, and the catalogue/source QFT discrepancy. Retained unresolved status and zero original-target progress; Optimality of the Fourier measurement ansatz over arbitrary quantum strategies is not established.

Coordinator checked Fourier-basis orthogonality, residue multiplicity d and the sign of the Hermitian quadratic-form matrix. Optimality over unrestricted measurements is not established. This additional check was model-assisted, not independent human verification.

\paragraph{Tools, process and contributions.}
Web browsing opened the primary arXiv abstract identified in the literature field; no full-paper proof or current-status search was performed. Python 3 was used for file preparation and structural checks; no numerical experiment or theorem prover was used. Contribution prompt A was followed by prompt B. GPT-6 Astra with ultra reasoning drafted the research and performed its stated self-audit; any additional coordinator contribution is identified in the verification and audit. A GPT-6-based Codex coordinator assembled the manuscript and checked repository integration; its exact variant and reasoning configuration were not exposed. No human mathematical contribution, independent human proof check or formal proof certification is supplied. Complete derivations and the consulted source references are included here; private reasoning transcripts are not published.

\paragraph{Final status and first remaining gap.}
\textbf{UNRESOLVED}. Optimality of the Fourier measurement ansatz over arbitrary quantum strategies is not established.

\section{Completion Estimate}
\noindent\textbf{COMPLETION: 0\%}
The calculation is an auxiliary or restricted attempt and establishes no newly verified portion of the full source target. Optimality of the Fourier measurement ansatz over arbitrary quantum strategies is not established.

\subsection{Sources}
\begin{itemize}
\item[\hypertarget{source:30000179:1}{S1}] ULAM.AI and the UnsolvedMath Contributors, UnsolvedMath, record 30000179 (OWR-785-005). Catalogue curation under CC BY 4.0; definitions and English exposition added locally.
\url{https://huggingface.co/datasets/ulamai/UnsolvedMath}.
\item[\hypertarget{source:30000179:2}{S2}] Entanglement and Decoherence: Mathematics and Physics of Quantum Information and Computation (2005), Oberwolfach Reports, DOI 10.4171/owr/2005/04, pp. 185--255; official report source and mathematical context.
\url{https://doi.org/10.4171/owr/2005/04}.
\item[\hypertarget{source:30000179:3}{S3}] Official EMS Press PDF of the same Oberwolfach report; PDF page 28, printed page 212.
\url{https://ems.press/content/serial-article-files/45978}.
\item[\hypertarget{source:30000179:4}{S4}] Collins et al., Bell inequalities for arbitrarily high-dimensional systems; reference in the saved dated catalogue literature review.
\url{https://arxiv.org/abs/quant-ph/0106024}.
\end{itemize}
\label{end:30000179}
\end{document}
